\corrigesChapitre{1}{Masse, volume et masse volumique}
\par\noindent\begin{minipage}{\linewidth}
\begin{corrige}{1}{Choisir l’instrument}
\begin{questions}
\item Une balance, après réglage du zéro ou tare du récipient.
\item Il s’agit d’un volume, et non d’une masse.
\end{questions}
\end{corrige}
\end{minipage}
\par\noindent\begin{minipage}{\linewidth}
\begin{corrige}{2}{Convertir sans confondre}
\begin{questions}
\item $0{,}250\times1\,000=250\unite{mL}$.
\item $35\unite{cm^3}=35\unite{mL}$.
\end{questions}
\end{corrige}
\end{minipage}
\par\noindent\begin{minipage}{\linewidth}
\begin{corrige}{3}{Lire une immersion}
\begin{questions}
\item $V=57-42=15\unite{mL}=15\unite{cm^3}$.
\item Le solide doit être entièrement immergé, insoluble et sans bulle adhérente.
\end{questions}
\par\smallskip\noindent\textbf{Pour comprendre.}\enspace On mesure le volume \emph{ajouté} par le solide : $V=V_f-V_i=57-42=15\unite{mL}=15\unite{cm^3}$. Vérifier immersion complète, absence de bulle et insolubilité.\par
\pcFigureEx{figures/cor-immersion.pdf}{Le volume du solide est la différence des deux lectures, pas le volume final seul.}
\end{corrige}
\end{minipage}
\par\noindent\begin{minipage}{\linewidth}
\begin{corrige}{4}{Une tare utile}
\begin{questions}
\item $m=132-82=50\unite g$.
\item Elle remettrait l’affichage à zéro avec le récipient vide ; la masse du liquide serait ensuite affichée directement.
\end{questions}
\end{corrige}
\end{minipage}
\par\noindent\begin{minipage}{\linewidth}
\begin{corrige}{5}{Calcul direct}
\begin{questions}
\item $\rho=32/40=0{,}80\unite{g/mL}$.
\item Non : $0{,}80<1{,}00$ dans les conditions indiquées.
\end{questions}
\par\smallskip\noindent\textbf{Pour comprendre.}\enspace Commencer par $\rho=m/V$ ; avec g et mL, on obtient g/mL. $32/40=0{,}80\unite{g/mL}$ ; comparer ensuite à $1{,}00\unite{g/mL}$ pour l’eau dans les conditions usuelles.\par
\end{corrige}
\end{minipage}
\par\noindent\begin{minipage}{\linewidth}
\begin{corrige}{6}{Une masse à prévoir}
\begin{questions}
\item $m=\rho V=2{,}7\times30=81\unite g$.
\item À conditions identiques, $162\unite g$ : la masse est proportionnelle au volume.
\end{questions}
\par\smallskip\noindent\textbf{Pour comprendre.}\enspace Isoler la masse : $m=\rho V=2{,}7\times30=81\unite g$. À même matériau et même température, doubler V double m ; le quotient m/V ne change pas.\par
\end{corrige}
\end{minipage}
\par\noindent\begin{minipage}{\linewidth}
\begin{corrige}{7}{Trouver le volume}
\begin{questions}
\item $V=m/\rho=79/7{,}9=10\unite{cm^3}$.
\item Cela correspond à $10\unite{mL}$.
\end{questions}
\par\smallskip\noindent\textbf{Pour comprendre.}\enspace Partir de $m=\rho V$, donc $V=m/\rho$. Les unités g et g/cm³ donnent cm³ : $79/7{,}9=10\unite{cm^3}=10\unite{mL}$.\par
\end{corrige}
\end{minipage}
\par\noindent\begin{minipage}{\linewidth}
\begin{corrige}{8}{Reconnaître une proportionnalité}
\begin{questions}
\item Oui : les trois quotients masse/volume valent $0{,}90$.
\item $\rho=0{,}90\unite{g/mL}$ ; le graphe est une droite passant par l’origine.
\end{questions}
\par\smallskip\noindent\textbf{Pour comprendre.}\enspace Calculer les trois quotients : $9/10=18/20=27/30=0{,}90$. Leur unité est g/mL. L’alignement sur une droite passant par l’origine est une autre lecture de la même proportionnalité.\par
\pcFigureEx{figures/cor-rho-points.pdf}{La droite passe par l’origine et les trois quotients sont identiques.}
\end{corrige}
\end{minipage}
\par\noindent\begin{minipage}{\linewidth}
\begin{corrige}{9}{Comparer à volume égal}
\begin{questions}
\item Aluminium : $2{,}70\times12=32{,}4\unite g$ ; fer : $7{,}9\times12=94{,}8\unite g$.
\item Le fer ; ici la comparaison des masses est pertinente parce que les volumes sont égaux.
\end{questions}
\end{corrige}
\end{minipage}
\par\noindent\begin{minipage}{\linewidth}
\begin{corrige}{10}{Changement d’unité}
\begin{questions}
\item $0{,}85\unite{g/mL}=850\unite{kg/m^3}$.
\item $2\unite L=2\,000\unite{mL}$ ; $m=0{,}85\times2\,000=1\,700\unite g=1{,}7\unite{kg}$.
\end{questions}
\par\smallskip\noindent\textbf{Pour comprendre.}\enspace $0{,}85\unite{g/mL}=850\unite{kg/m^3}$. Convertir aussi le volume : $2\unite L=2\,000\unite{mL}$, puis $m=0{,}85\times2\,000=1\,700\unite g=1{,}7\unite{kg}$.\par
\end{corrige}
\end{minipage}
\par\noindent\begin{minipage}{\linewidth}
\begin{corrige}{11}{Une bulle gênante}
\begin{questions}
\item Il est trop grand : le déplacement inclut le volume de la bulle.
\item À masse inchangée, diviser par un volume trop grand donne une masse volumique trop faible.
\end{questions}
\par\smallskip\noindent\textbf{Pour comprendre.}\enspace La masse mesurée séparément ne change pas. Si le dénominateur V est trop grand, le quotient m/V est trop petit : la masse volumique est sous-estimée.\par
\pcFigureEx{figures/cor-bulle.pdf}{La bulle augmente le volume apparent déplacé ; la masse du solide reste la même.}
\end{corrige}
\end{minipage}
\par\noindent\begin{minipage}{\linewidth}
\begin{corrige}{12}{Un objet creux}
\begin{questions}
\item $27/20=1{,}35\unite{g/cm^3}$.
\item Non : un objet creux peut avoir une masse volumique moyenne inférieure à celle du matériau de sa coque. Il faut examiner la présence d’une cavité.
\end{questions}
\pcFigureEx{figures/cor-coque.pdf}{Le volume extérieur comprend la cavité ; il diffère du volume de métal.}
\end{corrige}
\end{minipage}
\par\noindent\begin{minipage}{\linewidth}
\begin{corrige}{13}{Diagnostiquer un graphique}
\begin{questions}
\item Non en général : à volume nul, le récipient a encore une masse non nulle.
\item Soustraire la masse du récipient à toutes les masses, ou refaire les mesures après tare.
\end{questions}
\par\smallskip\noindent\textbf{Pour comprendre.}\enspace Le modèle d’une pesée sans tare est $m_{\text{mesurée}}=m_{\text{récipient}}+\rho V$. À V nul, l’ordonnée n’est donc pas nulle. Soustraire $m_{\text{récipient}}$ rétablit la masse du liquide.\par
\pcFigureEx{figures/cor-tare-graphe.pdf}{Soustraire la même masse de récipient à chaque mesure déplace toute la droite vers le bas.}
\end{corrige}
\end{minipage}
\par\noindent\begin{minipage}{\linewidth}
\begin{corrige}{14}{Comparer deux protocoles}
\begin{questions}
\item Dix billes donnent un déplacement plus grand ; une même difficulté de lecture représente alors une plus petite part du volume total.
\item $40/10=4{,}0\unite{mL}$ par bille. Vérifier qu’aucune bulle n’est emprisonnée.
\end{questions}
\par\smallskip\noindent\textbf{Pour comprendre.}\enspace Une même erreur de lecture représente une plus petite fraction d’un grand volume déplacé. Pour dix billes identiques : $V_{\text{bille}}=40/10=4\unite{mL}$. Vérifier que toutes sont immergées et sans bulle.\par
\end{corrige}
\end{minipage}
\par\noindent\begin{minipage}{\linewidth}
\begin{corrige}{15}{Un liquide qui se dilate}
\begin{questions}
\item Avant : $200/200=1{,}00\unite{g/mL}$ ; après : $200/204\approx0{,}980\unite{g/mL}$.
\item La même masse occupe plus de volume : la masse volumique diminue. La matière n’a pas disparu.
\end{questions}
\par\smallskip\noindent\textbf{Pour comprendre.}\enspace Avant : $\rho_1=200/200=1{,}00\unite{g/mL}$. Après : $\rho_2=200/204\approx0{,}980\unite{g/mL}$. La masse est conservée mais un volume plus grand diminue le quotient.\par
\end{corrige}
\end{minipage}
\par\noindent\begin{minipage}{\linewidth}
\phantomsection\label{cor-defi-1}\pcSolutionDefi{1}{La pièce mystérieuse}\begin{questions}
\item $\rho=43{,}2/16{,}0=2{,}70\unite{g/cm^3}$.
\item La mesure est compatible avec de l’aluminium compact dans les conditions de la référence ; elle ne constitue pas une preuve unique.
\item Recommencer la mesure après retrait des bulles et vérifier l’absence de cavité. On peut aussi améliorer la lecture du volume.
\end{questions}
\end{minipage}
\par\noindent\begin{minipage}{\linewidth}
\phantomsection\label{cor-activite-1}\pcSousTitre{Pistes pour l’investigation}
Pour chaque volume, soustraire la masse du récipient ou utiliser une tare correcte. Les points de l’eau doivent se situer près d’une droite de pente $1\unite{g/mL}$ vers 20 °C. Une concordance rend l’identification compatible, sans prouver une pureté parfaite. Un ménisque mal lu, une goutte perdue ou une température différente peut expliquer un écart.\par\end{minipage}
